\section{Effective-mass theory of shallow donors}\label{app:D.23}

\subsection*{The envelope equation}

The donor adds to the crystal Hamiltonian the potential $U(\vect{r})$ of its excess charge. Far from the core, at distances large compared with the lattice constant, it is the Coulomb potential of a unit charge, $-e^2/(4\pi r)$ in Heaviside--Lorentz units ($e^2 = 4\pi\alpha$). The host screens it by its static dielectric constant $\varepsilon$:
\begin{equation*}
  U(\vect{r}) = -\frac{e^2}{4\pi\varepsilon\,r} = -\frac{\alpha}{\varepsilon\,r} .
\end{equation*}
The dielectric constant is a pure number. For silicon, which being nonpolar has no lattice contribution, it is $\varepsilon = \varepsilon_0 = \varepsilon_\infty = 11.7$ (\cref{sec:5.5}), and from here on we write it $\varepsilon_0$, the static constant, as in \cref{sec:7.4}. Expand the donor state in the Bloch states of \cref{sec:3.2}, $\psi = \sum_{n\vect{k}}A_n(\vect{k})\,\psi_{n\vect{k}}$. The Schrödinger equation becomes
\begin{equation*}
  \bigl(\varepsilon_{n\vect{k}} - E\bigr)A_n(\vect{k}) + \sum_{n'\vect{k}'}\langle\psi_{n\vect{k}}|U|\psi_{n'\vect{k}'}\rangle\,A_{n'}(\vect{k}') = 0 .
\end{equation*}
A slowly varying $U$ simplifies the matrix element. Its Fourier components are concentrated at $|\vect{k} - \vect{k}'| \sim 1/a^*$, small compared with the zone, so for the cell-periodic parts $u_{n\vect{k}}$ of \cref{eq:3.2-bloch}, $\langle\psi_{n\vect{k}}|U|\psi_{n'\vect{k}'}\rangle \simeq \tilde U(\vect{k} - \vect{k}')\,\langle u_{n\vect{k}}|u_{n'\vect{k}'}\rangle_{\mathrm{cell}}$, where $\tilde U(\vect{p}) = \Vol^{-1}\int\dd^3r\,\ee^{-\ii\vect{p}\cdot\vect{r}}U(\vect{r})$ and the cell average is normalized to $\delta_{nn'}$ at $\vect{k} = \vect{k}'$. For $\vect{k}'$ close to $\vect{k}$ the overlap of different bands is of order $|\vect{k} - \vect{k}'|$ times a ratio of a velocity matrix element to an energy gap, by \cref{eq:3.2-dudk}, and it couples bands only at the next order in $1/(a^*k_{\mathrm{zone}})$. To leading order the bands decouple, and the conduction-band amplitude obeys
\begin{equation*}
  \bigl(\varepsilon_{c\vect{k}} - E\bigr)A(\vect{k}) + \sum_{\vect{k}'}\tilde U(\vect{k} - \vect{k}')\,A(\vect{k}') = 0 .
\end{equation*}

Near one conduction minimum $\vect{k}_j$, write $\vect{k} = \vect{k}_j + \delta\vect{k}$ and define the envelope $F_j(\vect{r}) = \Vol^{-1/2}\sum_{\delta\vect{k}}\ee^{\ii\delta\vect{k}\cdot\vect{r}}A(\vect{k}_j + \delta\vect{k})$. Multiplying by $\delta\vect{k}$ becomes $-\ii\nabla$ acting on $F_j$, and the convolution with $\tilde U$ becomes multiplication by $U(\vect{r})$. With the band near the minimum from \cref{sec:3.3}, $\varepsilon_{c,\vect{k}_j+\delta\vect{k}} = \Eg + \delta k_\parallel^2/2m^*_\parallel + \delta k_\perp^2/2m^*_\perp$ (energies from the valence maximum), and $z$ along the valley axis, the equation for the amplitude becomes the envelope equation \cref{eq:7.4-EMT}, with $E = \Eg - E_{\mathrm{b}}^{\mathrm{D}}$. The derivation assumed three things: that the envelope extends over many cells, $a^* \gg a$; that the binding is small compared with the energy over which the band departs from its quadratic form; and that the potential is the screened Coulomb potential everywhere that matters. The third fails within a cell or two of the core, which is the central-cell region treated below.

For $m^*_\parallel = m^*_\perp = m^*$, \cref{eq:7.4-EMT} is the hydrogen atom with $m_e \to m^*$ and $\alpha \to \alpha/\varepsilon_0$, and \cref{eq:7.4-hydrogenic} follows from the hydrogenic energies and length scale by the same substitution, with $a_0 = 1/(\alpha m_e) = \SI{0.529}{\angstrom}$ and $\mathrm{Ry} = \SI{13.606}{eV}$~\cite{mohr2025codata}.

\subsection*{Anisotropic valleys}

With unequal masses, \cref{eq:7.4-EMT} does not separate, and Kohn and Luttinger solved it variationally with an envelope squeezed along the heavy axis~\cite{kohn1955theory},
\begin{equation*}
  F(\vect{r}) = (\pi r_\perp^2r_\parallel)^{-1/2}\exp\Bigl[-\Bigl(\frac{x^2 + y^2}{r_\perp^2} + \frac{z^2}{r_\parallel^2}\Bigr)^{1/2}\Bigr] ,
\end{equation*}
with the radii $r_\perp$ across the valley axis and $r_\parallel$ along it, normalized to one. In the scaled coordinates $(x/r_\perp, y/r_\perp, z/r_\parallel)$ it is the hydrogen ground state $\ee^{-r}/\sqrt\pi$, for which each Cartesian component of $\langle|\nabla F|^2\rangle$ is $\frac13$. The kinetic energy is therefore $(2/3)(2m^*_\perp r_\perp^2)^{-1} + (1/3)(2m^*_\parallel r_\parallel^2)^{-1}$. For the potential, write $r = R\,(r_\perp^2\sin^2\theta + r_\parallel^2\cos^2\theta)^{1/2}$ in the scaled spherical coordinates $(R,\theta)$. The radial integral $\int_0^\infty R\,\ee^{-2R}\,\dd R = \frac14$ separates, and for $r_\parallel < r_\perp$
\begin{equation*}
  \Bigl\langle\frac1r\Bigr\rangle = \int_0^1\frac{\dd\mu}{\bigl[r_\perp^2 - (r_\perp^2 - r_\parallel^2)\mu^2\bigr]^{1/2}} = \frac{\arcsin e_{\mathrm{ecc}}}{r_\perp e_{\mathrm{ecc}}}, \qquad e_{\mathrm{ecc}} = \Bigl(1 - \frac{r_\parallel^2}{r_\perp^2}\Bigr)^{1/2} ,
\end{equation*}
Measuring lengths in units of $a^*_\perp = \varepsilon_0/(m^*_\perp\alpha)$ and energies in units of $\mathrm{Ry}^*_\perp = m^*_\perp\alpha^2/2\varepsilon_0^2$, with $e_{\mathrm{ecc}}$ the eccentricity of the spheroid, the energy is
\begin{equation}
  \frac{E(r_\perp,r_\parallel)}{\mathrm{Ry}^*_\perp} = \frac{2}{3r_\perp^2} + \frac{\gamma_m}{3r_\parallel^2} - \frac{2\arcsin e_{\mathrm{ecc}}}{r_\perp e_{\mathrm{ecc}}}, \qquad \gamma_m = \frac{m^*_\perp}{m^*_\parallel} .
  \label{eq:D.23-variational}
\end{equation}
For $\gamma_m = 1$ the minimum is at $r_\perp = r_\parallel = 1$ with $E = -1$, the hydrogenic result. For silicon, $\gamma_m = 0.191/0.916 = 0.209$, $\mathrm{Ry}^*_\perp = \SI{18.98}{meV}$ and $a^*_\perp = \SI{32.4}{\angstrom}$. Minimizing \cref{eq:D.23-variational} numerically gives $E = -1.564\,\mathrm{Ry}^*_\perp$ at $r_\perp = \SI{24.3}{\angstrom}$ and $r_\parallel = \SI{14.0}{\angstrom}$:
\begin{equation}
  E_{\mathrm{b}}^{\mathrm{D}}(\mathrm{Si}) = \SI{29.7}{meV} \qquad\text{(effective-mass theory, variational)} .
  \label{eq:D.23-Si}
\end{equation}
The radii are four and a half and two and a half lattice constants. The variational energy is an upper bound on the exact effective-mass energy.

\subsection*{Valleys and the central cell}

Silicon has six equivalent minima on the $\langle100\rangle$ axes. At the order of \cref{eq:7.4-EMT} each valley carries its own envelope, rotated with its axis, and the ground state is sixfold degenerate (twelvefold with spin): $\psi = \sum_{j=1}^6c_jF_j(\vect{r})\psi_{c\vect{k}_j}(\vect{r})$ with any coefficients $c_j$. The valleys couple through the part $\delta U$ of the donor potential that the envelope equation leaves out, which differs from $-\alpha/\varepsilon_0r$ only within about a lattice constant of the core, where the screening is incomplete and the ion has its own chemistry. Since $\delta U$ is short-ranged and each envelope is nearly constant over its range, the matrix elements between valley states are $\langle j|\delta U|j'\rangle \simeq F_j(0)F_{j'}(0)\int\dd^3r\,\psi^*_{c\vect{k}_j}\delta U\,\psi_{c\vect{k}_{j'}}$.

The tetrahedral symmetry of the site leaves three distinct elements: $\Delta_0$ on the diagonal, $\Delta_{\mathrm{g}}$ between opposite valleys and $\Delta_{\mathrm{f}}$ between perpendicular ones. The $6\times6$ matrix has the eigenvectors
\begin{equation}
  \begin{aligned}
  A_1:&\quad (1,1,1,1,1,1)/\sqrt6, & E &= \Delta_0 + \Delta_{\mathrm{g}} + 4\Delta_{\mathrm{f}} ,\\
  E:&\quad (1,1,-1,-1,0,0)/2 \text{ and } (1,1,1,1,-2,-2)/\sqrt{12}, & E &= \Delta_0 + \Delta_{\mathrm{g}} - 2\Delta_{\mathrm{f}} ,\\
  T_2:&\quad (1,-1,0,0,0,0)/\sqrt2 \text{ and its two partners}, & E &= \Delta_0 - \Delta_{\mathrm{g}} ,
  \end{aligned}
  \label{eq:D.23-vo}
\end{equation}
with the valleys ordered as $(+x,-x,+y,-y,+z,-z)$. In the $T_2$ states each pair of opposite valleys enters with opposite signs, and the four perpendicular couplings cancel. In the $A_1$ state all six valleys add in phase, so its density at the core is the largest, $|\psi(0)|^2 = 6|F(0)|^2|u_c(0)|^2$, and an attractive central-cell correction lowers it most. The $E$ and $T_2$ levels have nodes at the core and lie close to the effective-mass value.
